\documentclass{article}
\usepackage{apacite}
\usepackage{hyperref}
\usepackage{multicol}
\usepackage{natbib}
\usepackage{andi}
\usepackage[utf8]{inputenc}
\usepackage{array}
\usepackage[first=0,last=9]{lcg}
\setlength{\parindent}{0pt}
\usepackage{hyperref}
\usepackage{physics}
\usepackage{soul}   
\usepackage{tikz-cd}

\usepackage{courier} %% Sets font for listing as Courier.
\usepackage{listings, xcolor}
\lstset{
tabsize = 4, %% set tab space width
showstringspaces = false, %% prevent space marking in strings, string is defined as the text that is generally printed directly to the console
numbers = left, %% display line numbers on the left
commentstyle = \color{green}, %% set comment color
keywordstyle = \color{blue}, %% set keyword color
stringstyle = \color{red}, %% set string color
rulecolor = \color{black}, %% set frame color to avoid being affected by text color
basicstyle = \small \ttfamily , %% set listing font and size
breaklines = true, %% enable line breaking
numberstyle = \tiny,
}


\newcommand*{\rom}[1]{\expandafter\@slowromancap\romannumeral #1@}

\usetikzlibrary{calc,patterns,angles,quotes}
\usepackage{pgfplots}
\pgfplotsset{width=15cm,height=10cm,compat=1.9}

\pgfplotstableread{
x	y	delta-x	delta-y
-0.009	-6.3255	0.0284	0.0028
0.0611	-5.9201	0.0302	0.0019
0.2967	-5.6324	0.0375	0.0014
0.4689	-5.4092	0.0445	0.0011
0.4903	-5.2269	0.0455	0.0009
0.4822	-5.0728	0.0451	0.0008
0.5518	-4.9392	0.0485	0.0007
}{\mytable}

\begin{filecontents}{graphdata.dat}
x	y	delta-x	delta-y
-0.009	-6.3255	0.0284	0.0028
0.0611	-5.9201	0.0302	0.0019
0.2967	-5.6324	0.0375	0.0014
0.4689	-5.4092	0.0445	0.0011
0.4903	-5.2269	0.0455	0.0009
0.4822	-5.0728	0.0451	0.0008
0.5518	-4.9392	0.0485	0.0007
\end{filecontents}

\usepackage[english]{babel}
\usepackage{fancyhdr}
\usepackage{lastpage}
\pagestyle{fancy}
\fancyhf{}


\pagestyle{empty}% for cropping
\usepackage{tikz}
\usetikzlibrary{calc}
\usepackage{zref-savepos}
\usepackage{tabu}

\newcounter{NoTableEntry}
\renewcommand*{\theNoTableEntry}{NTE-\the\value{NoTableEntry}}

\newcommand*{\strike}[2]{%
  \multicolumn{1}{#1}{%
    \stepcounter{NoTableEntry}%
    \vadjust pre{\zsavepos{\theNoTableEntry t}}% top
    \vadjust{\zsavepos{\theNoTableEntry b}}% bottom
    \zsavepos{\theNoTableEntry l}% left
    \hspace{0pt plus 1filll}%
    #2% content
    \hspace{0pt plus 1filll}%
    \zsavepos{\theNoTableEntry r}% right
    \tikz[overlay]{%
      \draw
        let
          \n{llx}={\zposx{\theNoTableEntry l}sp-\zposx{\theNoTableEntry r}sp-\tabcolsep},
          \n{urx}={\tabcolsep},
          \n{lly}={\zposy{\theNoTableEntry b}sp-\zposy{\theNoTableEntry r}sp},
          \n{ury}={\zposy{\theNoTableEntry t}sp-\zposy{\theNoTableEntry r}sp}
        in
        (\n{llx}, \n{lly}) -- (\n{urx}, \n{ury})
      ;
    }% 
  }%
}

\usepackage{tikz}
\usetikzlibrary{calc}

\tikzset{
    right angle quadrant/.code={
        \pgfmathsetmacro\quadranta{{1,1,-1,-1}[#1-1]}     % Arrays for selecting quadrant
        \pgfmathsetmacro\quadrantb{{1,-1,-1,1}[#1-1]}},
    right angle quadrant=1, % Make sure it is set, even if not called explicitly
    right angle length/.code={\def\rightanglelength{#1}},   % Length of symbol
    right angle length=2ex, % Make sure it is set...
    right angle symbol/.style n args={3}{
        insert path={
            let \p0 = ($(#1)!(#3)!(#2)$) in     % Intersection
                let \p1 = ($(\p0)!\quadranta*\rightanglelength!(#3)$), % Point on base line
                \p2 = ($(\p0)!\quadrantb*\rightanglelength!(#2)$) in % Point on perpendicular line
                let \p3 = ($(\p1)+(\p2)-(\p0)$) in  % Corner point of symbol
            (\p1) -- (\p3) -- (\p2)
        }
    }
}
\pagestyle{fancy}
\fancyhf{}
\title{Using Programming to Solve Math Problems}
\author{Andi Tse}
\date{\today}
%\rhead{Instructor: Ms. Hung}
\lhead{Andi Tse}
\chead{Using Programming to Solve Math Problems}
\rfoot{BSSCC}
%\lfoot{Data Management}
\cfoot{Page \thepage \hspace{1pt} of \pageref{LastPage}}
\pagestyle{fancy}

 


\begin{document}

\maketitle
\tableofcontents
\newpage


\section{Combinatorics Problem}
\begin{custom-big}[Question 1]
Find the number of four digits integers with no repeating digits, and are divisible by five, and are greater than 4500. 

\end{custom-big}
\subsection*{Combinatorics Approach}
    \textbf{Case 1, the integer ends in 0:}\\
    \underline{Case 1.1, the integer starts with 4, the number of integers are:}\\
    \[\underline{1}\times \underline{5}\times \underline{7}\times \underline{1}=35\tag{The second digit has to be chosen from $5,6,7,8,9$}\]
    
    \underline{Case 1.2, the integer starts with 5 or above, the number of integers are:}\\
    \[\underline{5}\times \underline{8}\times \underline{7}\times \underline{1}=280\]
    
    
    \textbf{Case 2, the integer ends in 5:}\\
    \underline{Case 2.1, the integer starts with 4, the number of integers are:}\\
    \[\underline{1}\times \underline{4}\times \underline{7}\times \underline{1}=28\tag{The second digit has to be chosen from $6,7,8,9$}\]
    
    \underline{Case 2.2, the integer starts with 6 or above, the number of integers are:}\\
    \[\underline{4}\times \underline{8}\times \underline{7}\times \underline{1}=224\]

    Adding them all up, $35+280+28+224=567$\\
    Doing all these case works is very brain consuming. If allowed we can cheat with programming. 
    \pagebreak
\subsection*{Programming Approach}


    \begin{lstlisting}[language = Java , frame = trBL , firstnumber = last , escapeinside={(*@}{@*)}]
public class Main {

	public static void main(String[] args) throws IOException {
		int cnt = 0;
		for (int i = 1000; i < 100000; i++) {
			if (isValid(i))
				cnt++;
		}
		System.out.println(cnt);
	}

	public static boolean isValid(int n) {
		char[] digits = String.valueOf(n).toCharArray();
		Arrays.sort(digits);
		if (digits.length != 4)
			return false;
   // if it is not a 4 digits integer, it is invalid
		for (int i = 0; i < digits.length - 1; i++) {
			if (digits[i] == digits[i + 1])
				return false;
    // if the digits are repeated, the integer is invalid
		}

		if (n % 5 != 0)
			return false;
		if (n <= 4500)
			return false;
		return true;
	}
}

    
\end{lstlisting}
The program outputs 567, which matches in answer from Combinatorics Approach. After observing both methods, we can conclude that the question can be solved relatively easy with programming. 





\pagebreak

\section{Scientific Computation and the BigDecimal Data Type}
\begin{custom-big}[Question 3]
It is well known that if the square root of a natural number is not an integer, then it is irrational. The decimal expansion of such square roots is infinite without any repeating pattern at all.
The square root of two is $1.41421356237309504880 \cdots$, and the digital sum of the first one hundred decimal digits is 475.
For the first one hundred natural numbers, find the total of the digital sums of the first one hundred decimal digits for all the irrational square roots.
\end{custom-big}


\begin{lstlisting}
import java.math.*;

public class Main {
	public static void main(String[] args) {
		int a = 0;
		for (int i = 1; i <= 100; i++) {
			if (Math.sqrt(i) * Math.sqrt(i) == (int) Math.sqrt(i) * Math.sqrt(i))
				continue;
			a += sumSquareRootDigits100(i);
		}
		System.out.println(a);
	}

	public static int sumSquareRootDigits100(int n) {
		int a = 0;

		String d = new BigDecimal(n).sqrt(new MathContext(103)).toString();
		int cnt = 0;
		for (char c : d.toCharArray()) {
			if (c == '.')
				continue;
			cnt++;
			a += Character.getNumericValue(c);
			if (cnt == 100)
				break;
		}
		return a;
	}
} // end class
\end{lstlisting}
A further exploration can be found \href{https://en.wikipedia.org/wiki/Methods_of_computing_square_roots}{here}.


\pagebreak
\section{Data Management MDM4U1 Example (Game Fair)}
\begin{custom-big}[MDM4U1 Game Fair]
Given function $P(m)$, evaluate all $P(m)$ for all $m\in\mathbb{Z}$ such that $-200\leq m\leq 19800$
\[P\left(m\right)=\sum_{\substack{x=0\\ \operatorname{mod}(n+m-15x,100)=0}}^{\floor{\frac{n+m}{15}}}{n\choose x}\left(\frac{13}{64}\right)^{2x}{n-x\choose \frac{m+n-15x}{100}}\left(\frac{1}{128}\right)^{\frac{m+n-15x}{50}}\left(\frac{15707}{16384}\right)^{n-x-\frac{m+n-15x}{100}}\]


\end{custom-big}
Desmos does not support conditional statement for sigmas. To evaluate the summation we have to write computer programs or utilize softwares like Wolfram Alpha.
\medskip
\begin{lstlisting}[language = Java , frame = trBL , firstnumber = last , escapeinside={(*@}{@*)}]
// Author: Andi Tse
import java.util.*;
import java.io.*;
import java.lang.*;
import java.math.*;

public class GameFairCalculation {
	public static final int n = 200;

	public static void main(String[] args) throws IOException {
		File moneyFile = new File("Game Fair Probability Distribution Over Money.txt");
		PrintWriter fw = new PrintWriter(moneyFile);
		BigDecimal[] p = new BigDecimal[20002];

		for (int m = -200; m <= 19800; m++) {
			if (P(m).equals(new BigDecimal(0)))
				continue;
			p[m + 200] = P(m);
		}

		for (int m = -200; m <= 19800; m++) {
			if (p[m + 200] == null || p[m + 200].equals(new BigDecimal(0)))
				continue;
			fw.print("\\left(" + m + ", ");
			String s = String.valueOf(p[m + 200]);
			if (s.contains("E")) {
				s = s.substring(0, s.indexOf("E")) + "\\cdot 10^{" + s.substring(s.indexOf("E") + 1, s.length()) + "}";
			}
			fw.print(s);
			fw.println("\\right)");
		}
		fw.close();

		BigDecimal a = new BigDecimal(0);
		for (int m = -200; m < 0; m++) {
			if (p[m + 200] == null)
				continue;
			a = a.add(p[m + 200]);
		}
		System.out.println("P(m < 0) = " + a);
		a = new BigDecimal(0);
		for (int m = -200; m < 19800; m++) {
			if (p[m + 200] == null)
				continue;
			a = a.add(p[m + 200]);
		}
		System.out.println("P(all m) = " + a);

	}

	public static BigDecimal P(int x, int y) {
		return c(n, x).multiply(new BigDecimal(Math.pow(52.0 / 256, 2 * x)).multiply(c(n - x, y).multiply(
				new BigDecimal(Math.pow(1.0 / 128, 2 * y) * Math.pow(1 - Math.pow(52.0 / 256, 2)- 
						Math.pow(1.0 / 128,2),n - x - y)))));
	}

	public static BigDecimal P(int m) {

		BigDecimal a = new BigDecimal(0);
		for (int x = 0; x <= (n + m) / 15; x++) {
			if ((n + m - 15 * x) % 100 == 0) {
				a = a.add(c(n, x)
						.multiply(new BigDecimal(Math.pow(52.0 / 256, 2 * x)).multiply(c(n - x, (m + n - 15 * x) / 100)
								.multiply(new BigDecimal(Math.pow(1.0 / 128, (m + n - 15 * x) / 50.0)
										* Math.pow(1 - Math.pow(52.0 / 256, 2) - Math.pow(1.0 / 128, 2),
												n - x - (m + n - 15 * x) / 100.0))))));
			}
		}
		return a;
	}

	public static BigDecimal c(final int N, final int K) {
		// returns N choose K in BigDecimal
		BigInteger ret = BigInteger.ONE;
		for (int k = 0; k < K; k++) {
			ret = ret.multiply(BigInteger.valueOf(N - k)).divide(BigInteger.valueOf(k + 1));
		}
		return new BigDecimal(ret);
	}
}
\end{lstlisting}

The BigDecimal data type was chosen over the double data type for higher precision. The program calculates the probability of every possible profit $m$ by the player and stores in the txt file, which can be accessed here\footnote{\url{https://drive.google.com/file/d/1kBNWo8XWLyiEc9TjvVEXgateX26_UnDV/view?usp=share_link}}. The probability distribution graph can be accessed here\footnote{\url{https://www.desmos.com/calculator/nsgkph1xhd}}. The program outputs to the console
\begin{align*}
    P(m < 0) = 0.95288533845462508395
\end{align*}
Players have $95.29\%$ chance of losing money in Andi's game. 




\pagebreak
\section{Bidirectional Recurrence Problem}
\begin{custom-big}[Question 2]
Let:
$$
\begin{array}{ll}
x(0) & =0 \\
x(1) & =1 \\
x(2 k) & =\left(3 x(k)+2 x\left(\left\lfloor\frac{k}{2}\right\rfloor\right)\right) \bmod 2^{60} \text { for } k \geq 1, \text { where }\lfloor\rfloor \text { is the floor function } \\
x(2 k+1) & =\left(2 x(k)+3 x\left(\left\lfloor\frac{k}{2}\right\rfloor\right)\right) \bmod 2^{60} \text { for } k \geq 1 \\
y_n(k) & = \begin{cases}x(k) & \text { if } k \geq n \\
2^{60}-1-\max \left(y_n(2 k), y_n(2 k+1)\right) & \text { if } k<n\end{cases} \\
A(n) & =y_n(1)
\end{array}
$$
You are given:
$$
\begin{array}{ll}
x(2) & =3 \\
x(3) & =2 \\
x(4) & =11 \\
y_4(4) & =11 \\
y_4(3) & =2^{60}-9 \\
y_4(2) & =2^{60}-12 \\
y_4(1) & =A(4)=8 \\
A(10) & =2^{60}-34 \\
A\left(10^3\right) & =101881
\end{array}
$$
Find $A\left(10^{12}\right)$
\end{custom-big}



\subsection*{Solution}
    \begin{lstlisting}[language = Java , frame = trBL , firstnumber = last , escapeinside={(*@}{@*)}]
public class Main {
	public static final long mod = (long) Math.pow(2, 60);

	public static void main(String[] args) throws IOException {
		System.out.println(A((long) Math.pow(10, 12)));
	}

	public static long x(long n) {
		if (n == 0)
			return 0;
		if (n == 1)
			return 1;
		if (n == 2)
			return 3;
		if (n == 3)
			return 2;
		if (n == 4)
			return 11;
		if (n % 2 == 0) {
			long k = n / 2;
			return (3 * x(k) + 2 * x(k / 2)) % mod;
		} else {
			long k = (n - 1) / 2;
			return (2 * x(k) + 3 * x(k / 2)) % mod;

		}
	}

	public static long y(long n, long k) {
		if (n == 4 && k == 4)
			return 11;
		if (n == 4 && k == 3)
			return mod - 9;
		if (n == 4 && k == 2)
			return mod - 12;
		if (n == 4 && k == 1)
			return A(4);
		if (k >= n) {
			return x(k);
		} else {
			return mod - 1 - Math.max(y(n, 2 * k), y(n, 2 * k + 1));
		}
	}

	public static long A(long n) {
		if (n == 4)
			return 8;
		if (n == 10)
			return mod - 34;
		if (n == 1000)
			return 101881;
		return y(n, 1);
	}
}

\end{lstlisting}

\end{document}